O glaucoma é uma condição capaz de afetar os olhos, causando danos ao nervo ocular
e fazendo com que o paciente tenha seu campo visual reduzido aos poucos.
Se o glaucoma não for tratado, o paciente pode ser levado à cegueira. Entre os sintomas estão:
\begin{itemize}
    \item Perda gradual da visão lateral
    \item Visão embaçada
    \item Olhos vermelhos e inchados
    \item Sensibilidade a luz
\end{itemize}
A doença é bastante comum para pessoas acima de 40 anos, portanto, é recomendável uma visita anaual
ao médico oftamologista com o intuito de descartar a presença da doença e que esta se torne
intratável. Entre os exames realizados estão imagem do olho e o Tomografia de Coerência Óptica.

Dado isso, utilizando dados de imagens e clinícos, algoritmos de aprendizado de máquina, do
inglês \textit{Machine Learning}, vem sendo utilizados para dar suporte aos profissionais com
diagnósticos automatizados com rapidzez eficiência. Somado isso, é possível que os diagnósticos
tragam uma explicabilidade, dando ainda mais suporte e recursos que o profissional tenho uma decisão
mais assertiva quanto ao tratamento.

Diante do exposto, este projeto tem como intuito utilizar modelos de aprendizado de máquina para a
deteção do glaucoma. Além disso, o projeto também apresenta múltiplas possibilidades de modelagem do
problema, utilizando diferentes estratégias e seus benefícios. Ao final do projeto, não
somente é explicitado o desempenho do modelo como também sua eficiência computacional e a
explicabilidade por trás das respostas dos modelos de aprendizado de máquina.


